\chapter*{Lời Mở Đầu}
\addcontentsline{toc}{chapter}{Lời Mở Đầu}
\markboth{Lời Mở Đầu}{Lời Mở Đầu}

\begin{truyen}{Lời Mở Đầu}{Opening Words}
\chuhoa{C}{ó những câu học sinh nói ra khiến cả lớp cười ồ. Nhưng cười xong, người thầy vẫn phải nghĩ: phía sau câu lý sự đó là sự lười, sự sợ, hay một nỗi bối rối thật sự?}
\textit{(Some student excuses make the whole class laugh. But after the laughter, a teacher still wonders: behind that excuse, is it laziness, fear, or a real confusion?) }

Cuốn sách này gom lại 100 màn đối đáp rất ``đời lớp học'' giữa trò và thầy. Mỗi đoạn ngắn, mỗi câu nói quen, mỗi lý lẽ cùn đều xuất phát từ những tình huống rất thật: ngại học, sợ sai, sợ bị chê, không thấy ý nghĩa của việc cố gắng, hoặc chỉ đơn giản là muốn chống chế cho đỡ quê.
\textit{(This book gathers 100 very real classroom exchanges between students and teachers. Each short dialogue comes from familiar moments: not wanting to study, fearing mistakes, fearing judgment, not seeing the point of effort, or simply trying to avoid embarrassment.)}

Quyển này không viết để chê học sinh lý sự. Quyển này viết để cả thầy và trò cùng soi lại mình: trò hiểu vì sao người lớn cứ nhắc mãi, còn thầy nhớ rằng đằng sau sự bướng bỉnh đôi khi là một đứa trẻ đang chưa biết nói cho đúng nỗi lòng của mình.
\textit{(This volume is not written to mock students. It is written so both teachers and learners can look again at themselves: students may understand why adults keep repeating certain lessons, and teachers may remember that stubbornness sometimes hides a child who cannot yet express what is really inside.)}
\end{truyen}

\begin{baihoc}
Một câu chống chế nếu được lắng nghe đúng cách có thể trở thành một bài học rất sâu.
\textit{(An excuse, when listened to properly, can become a deep lesson.)}
\end{baihoc}
\ngancach